\section{Multiplexed QR Symbols}
\label{sec:symbol}

Both symbol designs studied here stack three binary QR layers, one per RGB channel, so
that each module shows one of eight colors: the color's red, green and blue bits are the
module's bits in the three layers. A layer bit of 1 is drawn with its channel off, which
matches the dark-module convention of a monochrome QR code~\cite{iso18004}.

\emph{HiQ}~\cite{yang2018hiq} encodes three independent messages, one standard QR code per
layer, so each layer can be decoded on its own once the colors are recovered.

\emph{qrism}, the design implemented in the open-source library this work builds on,
encodes one message. Its codeword stream is split into three equal parts, each written to
one channel with its own Reed-Solomon blocks, so a symbol decodes only when all three
channels do. The finder patterns use the same complementary color pairs as HiQ (green ring
and magenta stone top left, red and cyan top right, blue and yellow bottom left). Alignment
patterns, timing patterns, and format and version information are black and white, as in a
monochrome QR code.

\subsection{Decoding pipeline}

A decoder for either design runs four stages. (1)~\emph{Localization} finds the finder
and alignment patterns and fits a homography per region of the symbol. In the library used
here this runs on a min-channel projection of the photograph, which keeps colored finder
rings dark, with sub-pixel finder and alignment centers. (2)~\emph{Sampling} reads one pixel
at the projected center of every module. (3)~\emph{Color recovery} assigns each sample one
of the eight colors. (4)~\emph{Decoding} splits the color grid into its three layers and
runs standard Reed-Solomon decoding on each. This paper is about stage~(3); the experiments
in Section~\ref{sec:setup} hold stages (1), (2) and (4) fixed so that only color recovery
varies.
